\documentclass[UTF8,a4paper,zihao=-4,fontset=fandol]{ctexart}
\usepackage{wefjnu-exam}

\begin{document}
\wefjnutitle{2025--2026 学年第一学期}{《高等数学 C（上）》期末考试试题（A 卷）}
\wefjnumeta{适用专业　全校各相关专业}{适用年级　2025 级}
  {考试时间　2026 年 1 月 14 日 09:00}{任课教师　朱金才等}
\wefjnuinfo{闭卷 \quad 120 分钟 \quad 满分 100 分}

% 以下为示例。按原卷填写题型、题号、分值和试题。
\wefjnusection{一、选择题（每小题 3 分）}
\begin{enumerate}
  \item $x=0$ 是 $f(x)=\dfrac{x}{\tan x}$ 的\wefjnuoptionblank 间断点。
  \wefjnuoptions{可去}{跳跃}{无穷}{振荡}
\end{enumerate}

\wefjnusection{二、填空题（每小题 3 分）}
\begin{enumerate}
  \item $\displaystyle\lim_{x\to0}\frac{\sin x}{x}=$\wefjnuanswerblank。
\end{enumerate}

\wefjnuquestion{三}{8}
求极限 $\displaystyle\lim_{x\to1}\left(\frac{1}{x-1}-\frac{1}{\ln x}\right)$。
\end{document}
